\documentclass[a4paper]{beamer}

% \usepackage[utf8]{inputenc}
\usepackage[T1]{fontenc}
\usepackage[spanish]{babel}
\usepackage{amsmath, amssymb}
\usepackage{siunitx}
\usepackage{tikz}
\usepackage{wrapfig}
\usepackage{ragged2e}
\usetikzlibrary{calc,angles,quotes}
\usetikzlibrary{arrows.meta}
\usefonttheme[onlymath]{serif}
\setbeamertemplate{navigation symbols}{}

% ============================================================
% plano-inclinado.tex
% Macros TikZ para dibujar un plano inclinado y trabajar con
% puntos, tangente, normal, proyecciones y tramo rugoso.
%
% Uso en main.tex:
%   \usetikzlibrary{calc,angles,quotes}
%   % ============================================================
% plano-inclinado.tex
% Macros TikZ para dibujar un plano inclinado y trabajar con
% puntos, tangente, normal, proyecciones y tramo rugoso.
%
% Uso en main.tex:
%   \usetikzlibrary{calc,angles,quotes}
%   % ============================================================
% plano-inclinado.tex
% Macros TikZ para dibujar un plano inclinado y trabajar con
% puntos, tangente, normal, proyecciones y tramo rugoso.
%
% Uso en main.tex:
%   \usetikzlibrary{calc,angles,quotes}
%   \input{plano-inclinado.tex}
%
% Por ahora se administra UN plano por tikzpicture, porque las
% coordenadas internas usan nombres PI-... fijos.
% ============================================================

\makeatletter

% Variables internas para guardar las coordenadas y comparar
\newdimen\PI@xOne \newdimen\PI@yOne
\newdimen\PI@xTwo \newdimen\PI@yTwo
\newdimen\PI@xA   \newdimen\PI@yA
\newdimen\PI@xB   \newdimen\PI@yB

% Estado y valores por defecto de las opciones
\newif\ifPI@showprojection
\newif\ifPI@showfloor
\newif\ifPI@showangle

\pgfkeys{
  /plano/.is family,
  /plano,
  order/.store in=\PI@order,
  show projection/.is if=PI@showprojection,
  show projection/.default=true,
  show floor/.is if=PI@showfloor,
  show floor/.default=true,
  show angle/.is if=PI@showangle,
  show angle/.default=true,
  angle/.store in=\PI@anglelabel,
  angle radius/.store in=\PI@angleradius,
  projection style/.store in=\PI@projectionstyle,
  floor style/.store in=\PI@floorstyle,
  plane style/.store in=\PI@planestyle
}

% \planoInclinado[opciones]{punto1}{punto2}
\newcommand{\planoInclinado}[3][]{%
  % Restablecer todos los valores en cada llamada
  \PI@showprojectionfalse
  \PI@showfloorfalse
  \PI@showanglefalse
  \def\PI@order{bottomtop}%
  \def\PI@anglelabel{\theta}%
  \def\PI@angleradius{7mm}%
  \def\PI@projectionstyle{dashed}%
  \def\PI@floorstyle{dotted}%
  \def\PI@planestyle{thick}%
  \pgfkeys{/plano/.cd,#1}%
  
  % Coordenadas recibidas
  \coordinate (PI-one) at (#2);
  \coordinate (PI-two) at (#3);
  
  % Extraer coordenadas cartesianas para evaluar el sentido
  \pgfextractx{\PI@xOne}{\pgfpointanchor{PI-one}{center}}
  \pgfextracty{\PI@yOne}{\pgfpointanchor{PI-one}{center}}
  \pgfextractx{\PI@xTwo}{\pgfpointanchor{PI-two}{center}}
  \pgfextracty{\PI@yTwo}{\pgfpointanchor{PI-two}{center}}

  \def\PI@ordertopbottom{topbottom}%
  \def\PI@ordertopdown{topdown}%
  
  \ifx\PI@order\PI@ordertopbottom
    \coordinate (PI-A) at (PI-two);
    \coordinate (PI-B) at (PI-one);
    \PI@xA=\PI@xTwo \PI@xB=\PI@xOne
  \else\ifx\PI@order\PI@ordertopdown
    \coordinate (PI-A) at (PI-two);
    \coordinate (PI-B) at (PI-one);
    \PI@xA=\PI@xTwo \PI@xB=\PI@xOne
  \else
    \def\PI@orderbottomup{bottomup}%
    \ifx\PI@order\PI@orderbottomup
      \coordinate (PI-A) at (PI-one);
      \coordinate (PI-B) at (PI-two);
      \PI@xA=\PI@xOne \PI@xB=\PI@xTwo
    \else
      \coordinate (PI-A) at (PI-one);
      \coordinate (PI-B) at (PI-two);
      \PI@xA=\PI@xOne \PI@xB=\PI@xTwo
    \fi
  \fi\fi
  
  % Proyección horizontal de B sobre la altura de A
  \coordinate (PI-P) at (PI-B |- PI-A);
  
  % Geometría principal
  \draw[\PI@planestyle] (PI-A) -- (PI-B);
  \ifPI@showprojection
    \draw[\PI@projectionstyle] (PI-B) -- (PI-P);
  \fi
  \ifPI@showfloor
    \draw[\PI@floorstyle] (PI-A) -- (PI-P);
  \fi
  
  % Ángulo adaptado al sentido de la rampa
  \ifPI@showangle
    \ifdim\PI@xB>\PI@xA
      % La rampa asciende hacia la derecha
      \pic[
        draw,
        angle radius=\PI@angleradius,
        angle eccentricity=1.3,
        "$\PI@anglelabel$" {font=\small}
      ] {angle = PI-P--PI-A--PI-B};
    \else
      % La rampa asciende hacia la izquierda
      \pic[
        draw,
        angle radius=\PI@angleradius,
        angle eccentricity=1.3,
        "$\PI@anglelabel$" {font=\small}
      ] {angle = PI-B--PI-A--PI-P};
    \fi
  \fi
  
  % Guardar ángulo de la tangente A -> B y el de la normal superior.
  \pgfmathanglebetweenpoints{\pgfpointanchor{PI-A}{center}}{\pgfpointanchor{PI-B}{center}}%
  \edef\PI@tangentangle{\pgfmathresult}%
  \ifdim\PI@tangentangle pt>90pt
    \pgfmathsetmacro{\PI@normalangle}{\PI@tangentangle-90}%
  \else
    \pgfmathsetmacro{\PI@normalangle}{\PI@tangentangle+90}%
  \fi
  
  % Componentes unitarias disponibles para cálculos posteriores.
  \pgfmathsetmacro{\PI@tx}{cos(\PI@tangentangle)}%
  \pgfmathsetmacro{\PI@ty}{sin(\PI@tangentangle)}%
  \pgfmathsetmacro{\PI@nx}{cos(\PI@normalangle)}%
  \pgfmathsetmacro{\PI@ny}{sin(\PI@normalangle)}%
  \coordinate (PI-tangente) at (\PI@tx,\PI@ty);
  \coordinate (PI-normal) at (\PI@nx,\PI@ny);
}

% \puntoRampa{t}{nombre}
\newcommand{\puntoRampa}[2]{%
  \coordinate (#2) at ($(PI-A)!#1!(PI-B)$);%
}

% Atajo solicitado: calcula el punto y lo guarda como (PI-punto).
\newcommand{\atRampa}[1]{\puntoRampa{#1}{PI-punto}}

% \tangenteRampa{nombre} y \normalRampa{nombre}
\newcommand{\tangenteRampa}[1]{\coordinate (#1) at (\PI@tx,\PI@ty);}
\newcommand{\normalRampa}[1]{\coordinate (#1) at (\PI@nx,\PI@ny);}

% \proyeccionPiso{coordenada-origen}{nombre-destino}
\newcommand{\proyeccionPiso}[2]{\coordinate (#2) at (#1 |- PI-A);}

% \centroCirculoRampa{t}{radio}{nombre-centro}
\newcommand{\centroCirculoRampa}[3]{%
  \puntoRampa{#1}{PI-contacto-temp}%
  \coordinate (#3) at
    ($(PI-contacto-temp)+({(#2)*\PI@nx},{(#2)*\PI@ny})$);%
}

% \circuloRampa[estilo de dibujo]{t}{radio}{nombre-centro}
\newcommand{\circuloRampa}[4][thick]{%
    \centroCirculoRampa{#2}{#3}{#4}%
    \draw[#1] (#4) circle (#3);%
}

% \tramoRugoso[estilo]{t-inicial}{t-final}
\newcommand{\tramoRugoso}[3][very thick]{%
  \puntoRampa{#2}{PI-rugoso-inicio}%
  \puntoRampa{#3}{PI-rugoso-fin}%
  \draw[#1] (PI-rugoso-inicio) -- (PI-rugoso-fin);%
}

% \flechaTangente[estilo]{punto}{longitud-firmada}
\newcommand{\flechaTangente}[3][->,thick]{%
  \draw[#1] (#2) -- ++({(#3)*\PI@tx},{(#3)*\PI@ty});%
}

% \flechaNormal[estilo]{punto}{longitud-firmada}
\newcommand{\flechaNormal}[3][->,thick]{%
  \draw[#1] (#2) -- ++({(#3)*\PI@nx},{(#3)*\PI@ny});%
}

\makeatother

%
% Por ahora se administra UN plano por tikzpicture, porque las
% coordenadas internas usan nombres PI-... fijos.
% ============================================================

\makeatletter

% Variables internas para guardar las coordenadas y comparar
\newdimen\PI@xOne \newdimen\PI@yOne
\newdimen\PI@xTwo \newdimen\PI@yTwo
\newdimen\PI@xA   \newdimen\PI@yA
\newdimen\PI@xB   \newdimen\PI@yB

% Estado y valores por defecto de las opciones
\newif\ifPI@showprojection
\newif\ifPI@showfloor
\newif\ifPI@showangle

\pgfkeys{
  /plano/.is family,
  /plano,
  order/.store in=\PI@order,
  show projection/.is if=PI@showprojection,
  show projection/.default=true,
  show floor/.is if=PI@showfloor,
  show floor/.default=true,
  show angle/.is if=PI@showangle,
  show angle/.default=true,
  angle/.store in=\PI@anglelabel,
  angle radius/.store in=\PI@angleradius,
  projection style/.store in=\PI@projectionstyle,
  floor style/.store in=\PI@floorstyle,
  plane style/.store in=\PI@planestyle
}

% \planoInclinado[opciones]{punto1}{punto2}
\newcommand{\planoInclinado}[3][]{%
  % Restablecer todos los valores en cada llamada
  \PI@showprojectionfalse
  \PI@showfloorfalse
  \PI@showanglefalse
  \def\PI@order{bottomtop}%
  \def\PI@anglelabel{\theta}%
  \def\PI@angleradius{7mm}%
  \def\PI@projectionstyle{dashed}%
  \def\PI@floorstyle{dotted}%
  \def\PI@planestyle{thick}%
  \pgfkeys{/plano/.cd,#1}%
  
  % Coordenadas recibidas
  \coordinate (PI-one) at (#2);
  \coordinate (PI-two) at (#3);
  
  % Extraer coordenadas cartesianas para evaluar el sentido
  \pgfextractx{\PI@xOne}{\pgfpointanchor{PI-one}{center}}
  \pgfextracty{\PI@yOne}{\pgfpointanchor{PI-one}{center}}
  \pgfextractx{\PI@xTwo}{\pgfpointanchor{PI-two}{center}}
  \pgfextracty{\PI@yTwo}{\pgfpointanchor{PI-two}{center}}

  \def\PI@ordertopbottom{topbottom}%
  \def\PI@ordertopdown{topdown}%
  
  \ifx\PI@order\PI@ordertopbottom
    \coordinate (PI-A) at (PI-two);
    \coordinate (PI-B) at (PI-one);
    \PI@xA=\PI@xTwo \PI@xB=\PI@xOne
  \else\ifx\PI@order\PI@ordertopdown
    \coordinate (PI-A) at (PI-two);
    \coordinate (PI-B) at (PI-one);
    \PI@xA=\PI@xTwo \PI@xB=\PI@xOne
  \else
    \def\PI@orderbottomup{bottomup}%
    \ifx\PI@order\PI@orderbottomup
      \coordinate (PI-A) at (PI-one);
      \coordinate (PI-B) at (PI-two);
      \PI@xA=\PI@xOne \PI@xB=\PI@xTwo
    \else
      \coordinate (PI-A) at (PI-one);
      \coordinate (PI-B) at (PI-two);
      \PI@xA=\PI@xOne \PI@xB=\PI@xTwo
    \fi
  \fi\fi
  
  % Proyección horizontal de B sobre la altura de A
  \coordinate (PI-P) at (PI-B |- PI-A);
  
  % Geometría principal
  \draw[\PI@planestyle] (PI-A) -- (PI-B);
  \ifPI@showprojection
    \draw[\PI@projectionstyle] (PI-B) -- (PI-P);
  \fi
  \ifPI@showfloor
    \draw[\PI@floorstyle] (PI-A) -- (PI-P);
  \fi
  
  % Ángulo adaptado al sentido de la rampa
  \ifPI@showangle
    \ifdim\PI@xB>\PI@xA
      % La rampa asciende hacia la derecha
      \pic[
        draw,
        angle radius=\PI@angleradius,
        angle eccentricity=1.3,
        "$\PI@anglelabel$" {font=\small}
      ] {angle = PI-P--PI-A--PI-B};
    \else
      % La rampa asciende hacia la izquierda
      \pic[
        draw,
        angle radius=\PI@angleradius,
        angle eccentricity=1.3,
        "$\PI@anglelabel$" {font=\small}
      ] {angle = PI-B--PI-A--PI-P};
    \fi
  \fi
  
  % Guardar ángulo de la tangente A -> B y el de la normal superior.
  \pgfmathanglebetweenpoints{\pgfpointanchor{PI-A}{center}}{\pgfpointanchor{PI-B}{center}}%
  \edef\PI@tangentangle{\pgfmathresult}%
  \ifdim\PI@tangentangle pt>90pt
    \pgfmathsetmacro{\PI@normalangle}{\PI@tangentangle-90}%
  \else
    \pgfmathsetmacro{\PI@normalangle}{\PI@tangentangle+90}%
  \fi
  
  % Componentes unitarias disponibles para cálculos posteriores.
  \pgfmathsetmacro{\PI@tx}{cos(\PI@tangentangle)}%
  \pgfmathsetmacro{\PI@ty}{sin(\PI@tangentangle)}%
  \pgfmathsetmacro{\PI@nx}{cos(\PI@normalangle)}%
  \pgfmathsetmacro{\PI@ny}{sin(\PI@normalangle)}%
  \coordinate (PI-tangente) at (\PI@tx,\PI@ty);
  \coordinate (PI-normal) at (\PI@nx,\PI@ny);
}

% \puntoRampa{t}{nombre}
\newcommand{\puntoRampa}[2]{%
  \coordinate (#2) at ($(PI-A)!#1!(PI-B)$);%
}

% Atajo solicitado: calcula el punto y lo guarda como (PI-punto).
\newcommand{\atRampa}[1]{\puntoRampa{#1}{PI-punto}}

% \tangenteRampa{nombre} y \normalRampa{nombre}
\newcommand{\tangenteRampa}[1]{\coordinate (#1) at (\PI@tx,\PI@ty);}
\newcommand{\normalRampa}[1]{\coordinate (#1) at (\PI@nx,\PI@ny);}

% \proyeccionPiso{coordenada-origen}{nombre-destino}
\newcommand{\proyeccionPiso}[2]{\coordinate (#2) at (#1 |- PI-A);}

% \centroCirculoRampa{t}{radio}{nombre-centro}
\newcommand{\centroCirculoRampa}[3]{%
  \puntoRampa{#1}{PI-contacto-temp}%
  \coordinate (#3) at
    ($(PI-contacto-temp)+({(#2)*\PI@nx},{(#2)*\PI@ny})$);%
}

% \circuloRampa[estilo de dibujo]{t}{radio}{nombre-centro}
\newcommand{\circuloRampa}[4][thick]{%
    \centroCirculoRampa{#2}{#3}{#4}%
    \draw[#1] (#4) circle (#3);%
}

% \tramoRugoso[estilo]{t-inicial}{t-final}
\newcommand{\tramoRugoso}[3][very thick]{%
  \puntoRampa{#2}{PI-rugoso-inicio}%
  \puntoRampa{#3}{PI-rugoso-fin}%
  \draw[#1] (PI-rugoso-inicio) -- (PI-rugoso-fin);%
}

% \flechaTangente[estilo]{punto}{longitud-firmada}
\newcommand{\flechaTangente}[3][->,thick]{%
  \draw[#1] (#2) -- ++({(#3)*\PI@tx},{(#3)*\PI@ty});%
}

% \flechaNormal[estilo]{punto}{longitud-firmada}
\newcommand{\flechaNormal}[3][->,thick]{%
  \draw[#1] (#2) -- ++({(#3)*\PI@nx},{(#3)*\PI@ny});%
}

\makeatother

%
% Por ahora se administra UN plano por tikzpicture, porque las
% coordenadas internas usan nombres PI-... fijos.
% ============================================================

\makeatletter

% Variables internas para guardar las coordenadas y comparar
\newdimen\PI@xOne \newdimen\PI@yOne
\newdimen\PI@xTwo \newdimen\PI@yTwo
\newdimen\PI@xA   \newdimen\PI@yA
\newdimen\PI@xB   \newdimen\PI@yB

% Estado y valores por defecto de las opciones
\newif\ifPI@showprojection
\newif\ifPI@showfloor
\newif\ifPI@showangle

\pgfkeys{
  /plano/.is family,
  /plano,
  order/.store in=\PI@order,
  show projection/.is if=PI@showprojection,
  show projection/.default=true,
  show floor/.is if=PI@showfloor,
  show floor/.default=true,
  show angle/.is if=PI@showangle,
  show angle/.default=true,
  angle/.store in=\PI@anglelabel,
  angle radius/.store in=\PI@angleradius,
  projection style/.store in=\PI@projectionstyle,
  floor style/.store in=\PI@floorstyle,
  plane style/.store in=\PI@planestyle
}

% \planoInclinado[opciones]{punto1}{punto2}
\newcommand{\planoInclinado}[3][]{%
  % Restablecer todos los valores en cada llamada
  \PI@showprojectionfalse
  \PI@showfloorfalse
  \PI@showanglefalse
  \def\PI@order{bottomtop}%
  \def\PI@anglelabel{\theta}%
  \def\PI@angleradius{7mm}%
  \def\PI@projectionstyle{dashed}%
  \def\PI@floorstyle{dotted}%
  \def\PI@planestyle{thick}%
  \pgfkeys{/plano/.cd,#1}%
  
  % Coordenadas recibidas
  \coordinate (PI-one) at (#2);
  \coordinate (PI-two) at (#3);
  
  % Extraer coordenadas cartesianas para evaluar el sentido
  \pgfextractx{\PI@xOne}{\pgfpointanchor{PI-one}{center}}
  \pgfextracty{\PI@yOne}{\pgfpointanchor{PI-one}{center}}
  \pgfextractx{\PI@xTwo}{\pgfpointanchor{PI-two}{center}}
  \pgfextracty{\PI@yTwo}{\pgfpointanchor{PI-two}{center}}

  \def\PI@ordertopbottom{topbottom}%
  \def\PI@ordertopdown{topdown}%
  
  \ifx\PI@order\PI@ordertopbottom
    \coordinate (PI-A) at (PI-two);
    \coordinate (PI-B) at (PI-one);
    \PI@xA=\PI@xTwo \PI@xB=\PI@xOne
  \else\ifx\PI@order\PI@ordertopdown
    \coordinate (PI-A) at (PI-two);
    \coordinate (PI-B) at (PI-one);
    \PI@xA=\PI@xTwo \PI@xB=\PI@xOne
  \else
    \def\PI@orderbottomup{bottomup}%
    \ifx\PI@order\PI@orderbottomup
      \coordinate (PI-A) at (PI-one);
      \coordinate (PI-B) at (PI-two);
      \PI@xA=\PI@xOne \PI@xB=\PI@xTwo
    \else
      \coordinate (PI-A) at (PI-one);
      \coordinate (PI-B) at (PI-two);
      \PI@xA=\PI@xOne \PI@xB=\PI@xTwo
    \fi
  \fi\fi
  
  % Proyección horizontal de B sobre la altura de A
  \coordinate (PI-P) at (PI-B |- PI-A);
  
  % Geometría principal
  \draw[\PI@planestyle] (PI-A) -- (PI-B);
  \ifPI@showprojection
    \draw[\PI@projectionstyle] (PI-B) -- (PI-P);
  \fi
  \ifPI@showfloor
    \draw[\PI@floorstyle] (PI-A) -- (PI-P);
  \fi
  
  % Ángulo adaptado al sentido de la rampa
  \ifPI@showangle
    \ifdim\PI@xB>\PI@xA
      % La rampa asciende hacia la derecha
      \pic[
        draw,
        angle radius=\PI@angleradius,
        angle eccentricity=1.3,
        "$\PI@anglelabel$" {font=\small}
      ] {angle = PI-P--PI-A--PI-B};
    \else
      % La rampa asciende hacia la izquierda
      \pic[
        draw,
        angle radius=\PI@angleradius,
        angle eccentricity=1.3,
        "$\PI@anglelabel$" {font=\small}
      ] {angle = PI-B--PI-A--PI-P};
    \fi
  \fi
  
  % Guardar ángulo de la tangente A -> B y el de la normal superior.
  \pgfmathanglebetweenpoints{\pgfpointanchor{PI-A}{center}}{\pgfpointanchor{PI-B}{center}}%
  \edef\PI@tangentangle{\pgfmathresult}%
  \ifdim\PI@tangentangle pt>90pt
    \pgfmathsetmacro{\PI@normalangle}{\PI@tangentangle-90}%
  \else
    \pgfmathsetmacro{\PI@normalangle}{\PI@tangentangle+90}%
  \fi
  
  % Componentes unitarias disponibles para cálculos posteriores.
  \pgfmathsetmacro{\PI@tx}{cos(\PI@tangentangle)}%
  \pgfmathsetmacro{\PI@ty}{sin(\PI@tangentangle)}%
  \pgfmathsetmacro{\PI@nx}{cos(\PI@normalangle)}%
  \pgfmathsetmacro{\PI@ny}{sin(\PI@normalangle)}%
  \coordinate (PI-tangente) at (\PI@tx,\PI@ty);
  \coordinate (PI-normal) at (\PI@nx,\PI@ny);
}

% \puntoRampa{t}{nombre}
\newcommand{\puntoRampa}[2]{%
  \coordinate (#2) at ($(PI-A)!#1!(PI-B)$);%
}

% Atajo solicitado: calcula el punto y lo guarda como (PI-punto).
\newcommand{\atRampa}[1]{\puntoRampa{#1}{PI-punto}}

% \tangenteRampa{nombre} y \normalRampa{nombre}
\newcommand{\tangenteRampa}[1]{\coordinate (#1) at (\PI@tx,\PI@ty);}
\newcommand{\normalRampa}[1]{\coordinate (#1) at (\PI@nx,\PI@ny);}

% \proyeccionPiso{coordenada-origen}{nombre-destino}
\newcommand{\proyeccionPiso}[2]{\coordinate (#2) at (#1 |- PI-A);}

% \centroCirculoRampa{t}{radio}{nombre-centro}
\newcommand{\centroCirculoRampa}[3]{%
  \puntoRampa{#1}{PI-contacto-temp}%
  \coordinate (#3) at
    ($(PI-contacto-temp)+({(#2)*\PI@nx},{(#2)*\PI@ny})$);%
}

% \circuloRampa[estilo de dibujo]{t}{radio}{nombre-centro}
\newcommand{\circuloRampa}[4][thick]{%
    \centroCirculoRampa{#2}{#3}{#4}%
    \draw[#1] (#4) circle (#3);%
}

% \tramoRugoso[estilo]{t-inicial}{t-final}
\newcommand{\tramoRugoso}[3][very thick]{%
  \puntoRampa{#2}{PI-rugoso-inicio}%
  \puntoRampa{#3}{PI-rugoso-fin}%
  \draw[#1] (PI-rugoso-inicio) -- (PI-rugoso-fin);%
}

% \flechaTangente[estilo]{punto}{longitud-firmada}
\newcommand{\flechaTangente}[3][->,thick]{%
  \draw[#1] (#2) -- ++({(#3)*\PI@tx},{(#3)*\PI@ty});%
}

% \flechaNormal[estilo]{punto}{longitud-firmada}
\newcommand{\flechaNormal}[3][->,thick]{%
  \draw[#1] (#2) -- ++({(#3)*\PI@nx},{(#3)*\PI@ny});%
}

\makeatother

\title[About Beamer] %optional
{Física de los Sist. de Partículas - CB020}

\subtitle{Resolución de ejercicio de cuerpo rígido en plano inclinado}

\author[G5] % (optional, for multiple authors)
{
	Grupo 5
}

\titlegraphic{\includegraphics[width=1\textwidth]{logofiuba.png}}

\institute[VFU] % (optional)
{
	Facultad de Ingeniería\\
	Universidad de Buenos Aires
}

\date[VLC 2021] % (optional)
{Exposición Grupal, Octubre 2026}


% \logo{\includegraphics[width=1\textwidth]{logofiuba.png}}



\newcommand{\tcm}{\ensuremath\text{\tiny CM}}
\newcommand{\tcir}{\ensuremath\text{\tiny CIR}}

\newcommand{\fr}{\ensuremath{\mathit{f\kern-0.12em r}}}
\newcommand{\fre}{\ensuremath{\mathit{f\kern-0.05em r}\kern-0.05em_{e}}}

\newcommand{\uveci}{\hat{\imath}}
\newcommand{\uvecj}{\hat{\jmath}}

\newcommand{\syscoord}[2]{
      \begin{scope}[rotate around={#2:(#1)}]
        
        % Ejes y elementos vectoriales
        \draw[-latex] (#1) -- ++(-.7,0) node[above right] {$x$};
        \draw[-latex] (#1) -- ++(0,-.7) node[right] {$y$};
        
        \draw[-stealth] (#1) ++(-0.45,-0.1) arc (0:90:-0.35) 
              node[midway,left] {$\vec{\tau}$};
        
        \draw (#1) circle(.1);
        \filldraw (#1) circle(1pt) node[left] {$z$};

        % Etiqueta del origen "O"
        \node[right] at (#1) {O};

      \end{scope}
}


\makeatletter

% ============================================================
% cilindro.tex
%
% Uso:
%   \dibujarCilindro[opciones]{radio}
%
% Opciones:
%   rozamiento=0|1       0: dibuja fr_e; 1: la oculta
%   coordenadas=true|false
%   ejes=true|false
%   origen={(x,y)}
%
% Requiere:
%   TikZ y la macro \syscoord definida en main.tex
% ============================================================

\pgfkeys{
    /cilindro/.is family,
    /cilindro,
    rozamiento/.store in=\CI@rozamiento,
    coordenadas/.store in=\CI@coordenadas,
    ejes/.store in=\CI@ejes,
    origen/.store in=\CI@origen
}

\newcommand{\dibujarCilindro}[2][]{%
    \begingroup

    % Valores predeterminados
    \def\CI@rozamiento{0}%
    \def\CI@coordenadas{false}%
    \def\CI@ejes{false}%
    \def\CI@origen{0,0}%

    % Leer opciones
    \pgfkeys{/cilindro/.cd,#1}%

    \def\R{#2}%

    % Centro del cilindro
    \coordinate (O) at (\CI@origen);

    % Círculo y centro de masa
    \draw (O) circle (\R);
    \fill (O) circle (0.05) node[right]{\tiny CM};

    % Fuerzas en el centro
    \draw[-Stealth,thick,olive]
        (O) -- ++({-0.5*\R},0)
        node[above]{$\vec{F}$};

    \draw[-Stealth,thick]
        (O) -- ++(0,{-0.5*\R})
        node[right]{$m\vec{g}$};

    % Referencias giradas respecto al centro
    \begin{scope}[rotate around={-30:(O)}]

        % Sistema de coordenadas opcional
        \def\CI@true{true}%
        \ifx\CI@coordenadas\CI@true
            \pgfmathsetmacro{\CI@coordx}{1.25*\R}%
            \pgfmathsetmacro{\CI@coordy}{1.25*\R}%
            \coordinate (CI-syscoord)
                at ($(O)+(\CI@coordx,\CI@coordy)$);
            \syscoord{CI-syscoord}{0}%
        \fi

        % Ejes de referencia opcionales
        \ifx\CI@ejes\CI@true
            \draw[dashed]
                ($(O)+({-1.25*\R},0)$)
                -- ($(O)+({1.25*\R},0)$);

            \draw[dashed]
                ($(O)+(0,{-1.25*\R})$)
                -- ($(O)+(0,{1.25*\R})$);

            \draw
                ($(O)+(0,{-\R})$)
                -- ++({-0.75*\R},0)
                -- ++({1.5*\R},0);
        \fi

        % Fuerza normal
        \draw[-Stealth,thick,blue]
            ($(O)+(0,{-\R})$)
            -- ++(0,{0.5*\R})
            node[midway,left]{$\vec{N}$};

        % Fuerza de rozamiento
        \def\CI@cero{0}%
        \ifx\CI@rozamiento\CI@cero
            \draw[-Stealth,thick,red]
                ($(O)+(0,{-\R})$)
                -- ++({0.5*\R},0)
                node[midway,below]{$\vec{\fr}_e$};
        \fi

    \end{scope}

    \endgroup
}

\makeatother


\begin{document}


\frame{\titlepage}
\begin{frame}
	% \justifying

	\frametitle{Situación}
	Un cilindro homogéneo de radio $R = \qty{0.10}{\meter}$ y masa $M = \qty{2}{\kilo\gram}$ asciende rodando sin deslizar, por un plano inclinado un ángulo $\alpha = \qty{30}{\degree}$ con la horizontal.
	\vspace{-2em}
	\begin{wrapfigure}{r}{0.4\textwidth}
		% \centering
		\begin{tikzpicture}
			\def\R{.4}
			\planoInclinado[
				show angle,
				show projection,
				show floor,
				angle=\theta,
				projection style=dashed,
				floor style=thick,
			]{4,0}{0,2}

			\tramoRugoso[color=red, very thick]{.1}{.7}
			\centroCirculoRampa{.4}{\R}{C1}
			\draw (C1) circle (\R);
			\atRampa{.4}
			\node [below] at (PI-punto) {A};

			\centroCirculoRampa{.7}{\R}{C1}
			\draw (C1) circle (\R);
			\atRampa{.7}
			\node [below] at (PI-punto) {B};
		\end{tikzpicture}
		% \caption{}
		\label{fig:}
	\end{wrapfigure}
	\begin{justify}
		En todo momento se le aplica en el CM una fuerza horizontal, con el sentido indicado en la figura, de módulo constante $F = \qty{20}{\newton}$.
		El cuerpo parte del reposo en la posición A, en el que la altura del CM es de $h_{\tcm_A}$ = \qty{1}{\meter}, respecto a la superficie horizontal indicada.
		En el punto B el CM alcanza una altura de $\qty{3}{\meter}$, y a partir de ese punto el cuerpo se mueve en un tramo liso (el rozamiento puede considerarse despreciable).
	\end{justify}

\end{frame}

\begin{frame}
	\frametitle{Problema}
	Se pide
	\begin{itemize}
		\item Para cada tramo de la trayectoria, realizar el DCL para el cilindro, plantear las ecuaciones dinámicas de movimiento, indicando claramente el sistema de referencia y el sistema de coordenadas elegidos.
		\item Calcular la aceleración angular y la del centro de masa del cilindro durante la subida en el tramo $\overline{AB}$ y la aceleración angular en el tramo liso.
		\item Utilizando consideraciones energéticas, encontrar la velocidad del centro de masa del cilindro en un punto C de la trayectoria en el que la altura del centro de masa es de \qty{4}{\meter} respecto del plano horizontal indicado.
	\end{itemize}


\end{frame}


%% Descomposicion tramo B
\begin{frame}
	\frametitle{DCL y Ecuaciones Dinámicas: (Tramo B)}
	\begin{figure}
		% \centering
		\begin{tikzpicture}
			\def\R{1.5}
			\coordinate (O) at (0,0) ;


			\draw (O) circle (\R);
			\fill (O) circle(.05) node[right]{\tiny CM};

			\draw[-Stealth, thick, olive] (O) -- ++ (-0.5 * \R,0) node [left]{$\vec{F}$};
			\draw[-Stealth, thick] (O) -- ++ (0, -0.5 * \R) node [right]{$m\vec{g}$};
			\begin{scope}[rotate around={-30:(O)}]
				\syscoord{2.5,2.25}{0}

				\draw [dashed] (O) ++ (-1.25 * \R,0) -- ++ (2 * 1.25 * \R,0);
				\draw [dashed] (O) ++ (0,-1.25 * \R) -- ++ (0,2 * 1.25 * \R);
				\draw (O) ++ (0,-\R) -- ++(-0.75*\R,0) -- ++ (1.5*\R,0);

				\draw[-Stealth, thick, blue] (O) ++ (0,-\R) -- ++(0,0.5*\R) node[midway, left]{$\vec{N}$};
			\end{scope}
		\end{tikzpicture}
		% \caption{}
		% \label{fig:}
	\end{figure}
	$$ \begin{array}{l c r}

			\sum{F_y} = m\,a_y &
			\sum{F_x} = m\,a_x &
			\sum{\vec{\tau}_{\text{\tiny CM}}} = \sum{\vec{r}_i \times \vec{F_i}} = I_{\text{\tiny CM}}\,\vec{\gamma}
		\end{array} $$

	$$ F\,\sin(\alpha) + m\,g\,\cos(\alpha) - N = 0 $$
	$$ F\,\cos(\alpha) - m\,g\,\sin(\alpha) = m\,a_x $$
	$$ \tau_c = \vec{0} \times m\,\vec{g} + \vec{0} \times \vec{F} + R\,\uvecj \times \vec{N} $$
\end{frame}


%%%% Punto B

\begin{frame}
	\frametitle{Cálculo de aceleración angular (Tramo B)}

	Se observa en el DCL que todas las fuerzas tienen torque nulo, ya que están aplicadas sobre el CM o son paralelas a su vector posición.
	% \vspace{.35cm}
	\begin{columns}
		\column{0.3 \textwidth}
		\begin{figure}
			\begin{tikzpicture}
				\dibujarCilindro[rozamiento=1]{1.5}
			\end{tikzpicture}
		\end{figure}
		\column{0.5 \textwidth}
		$$ \tau_c = \vec{0} \times m\,\vec{g} + \vec{0} \times \vec{F} + R\,\uvecj \times \vec{N} $$
		\begin{equation} \label{eq:torque_conservation}
			\sum{\vec{\tau}} = \frac{\mathrm{d}\vec{L}}{\mathrm{d}t} = I_{\text{\tiny CM}} \,\vec{\gamma} = \vec{0}
		\end{equation}
		$$ \implies \gamma = 0 $$
	\end{columns}
	\vspace{.5cm}
	Esto nos indica que a partir de este punto $\omega$ es constante.

	\vspace{.4cm}
	Si en cambio el cilindro asciende rodando sin deslizar, \textquestiondown hacia dónde debería apuntar $\vec{\fr}_e$?
\end{frame}

%% Descomposicion tramo A
\begin{frame}
	\frametitle{DCL y Ecuaciones Dinámicas: (Tramo A)}

	\begin{columns}
		\column{0.6\textwidth}
		\begin{figure}
			% \centering
			\begin{tikzpicture}
				\def\R{1.5}
				\coordinate (O) at (0,0) ;


				\draw (O) circle (\R);
				\fill (O) circle(.05) node[right]{\tiny CM};

				\draw[-Stealth, thick, olive] (O) -- ++ (-0.5 * \R,0) node [left]{$\vec{F}$};
				\draw[-Stealth, thick] (O) -- ++ (0, -0.5 * \R) node [right]{$m\vec{g}$};
				\begin{scope}[rotate around={-30:(O)}]
					\syscoord{2,2}{0}

					\draw [dashed] (O) ++ (-1.25 * \R,0) -- ++ (2 * 1.25 * \R,0);
					\draw [dashed] (O) ++ (0,-1.25 * \R) -- ++ (0,2 * 1.25 * \R);
					\draw (O) ++ (0,-\R) -- ++(-0.75*\R,0) -- ++ (1.5*\R,0);

					\draw[-Stealth, thick, blue] (O) ++ (0,-\R) -- ++(0,0.5*\R) node[midway, left]{$\vec{N}$};
					\draw[-Stealth, thick, red] (O) ++ (0,-\R) -- ++(0.5*\R,0) node[midway, above right]{$\vec{\fr}_e$};
				\end{scope}
			\end{tikzpicture}
		\end{figure}

		\column{0.5\textwidth}
		$$\sum{F_y} = m\,a_y $$
		$$\sum{F_x} = m\,a_x $$
		$$\sum{\vec{\tau}_{\text{\tiny CM}}} = \sum{\vec{r}_i \times \vec{F_i}} = I_{\text{\tiny CM}}\,\vec{\gamma} $$
	\end{columns}

	$$ F\,\sin(\alpha) + m\,g\,\cos(\alpha) - N = 0 $$
	$$ F\,\cos(\alpha) - m\,g\,\sin(\alpha) - \fre = m\,a_x $$
	$$ \tau_c = \vec{0} \times m\,\vec{g} + \vec{0} \times \vec{F} + R\,\uvecj \times \vec{N} + R\,\uvecj \times \vec{\fre} $$

	\vspace{1em}
	Esta vez, con la contribución de $\vec{\fr}_e$, el torque deja de ser nulo.
\end{frame}


\begin{frame}
	\frametitle{Campo de aceleraciones durante RSD (Tramo A)}
	Dada la condición de rodadura sin deslizar, el CIR solo tiene aceleración normal, y obtenemos la aceleración del CM como
	$$ \vec{a}_{\tcir} = \vec{\omega} \times (\vec{\omega} \times -\vec{r}_{\tcm/\tcir}) $$
	Planteamos campo de aceleraciones
	$$ \vec{a}_{\tcm} = \vec{a}_\tcir + \vec{a}_{\tcm/\tcir} = \vec{a}_{\tcir} + \vec{\omega} \times (\vec{\omega} \times \vec{r}_{\tcm/\tcir}) + \vec{\gamma} \times \vec{r}_{\tcm/\tcir} $$
	% $$ \vec{a}_{\text{\tiny CM}} = \vec{r}_{\text{\tiny CM}/\text{\tiny CIR}} \times \vec{\gamma} \implies a_{\text{\tiny CM}_x} = \gamma\,R $$
	$$ \vec{a}_{\text{\tiny CM}} = \vec{\gamma} \times \vec{r}_{\tcm/\tcir} $$
	$$ = \left|\begin{matrix}
			\uveci & \uvecj & \hat{k} \\
			0      & 0      & \gamma  \\
			0      & -R     & 0
		\end{matrix}\right| $$
	$$ = \gamma\,R $$

\end{frame}

\begin{frame}
	\frametitle{Cálculo de aceleración angular y del CM (Tramo A)}
	\begin{figure}
		\begin{tikzpicture}
			\dibujarCilindro{1.5}
		\end{tikzpicture}
	\end{figure}
	De las ecuaciones dinámicas, despejamos $\fre$
	% $$ I_{\text{\tiny CM}}\,\vec{\gamma} = R\,\uvecj \times -\fre \, \uveci = \fre\,R\,\hat{k} = \frac{1}{2}\,m\,R^2 \gamma\,\hat{k} \implies \fre = \frac{1}{2}\,m\,R\,\gamma$$
	% \vspace{.15em}
	$$ I_{\text{\tiny CM}}\,\vec{\gamma} = R\,\uvecj \times -\fre \, \uveci$$ $$ \implies  \fre\,R\,\hat{k} = \frac{1}{2}\,m\,R^2 \gamma\,\hat{k} $$
	$$ \implies \fre = \frac{1}{2}\,m\,R\,\gamma$$
	$$ \fre = F\,\cos(\alpha) - m\,g\,\sin(\alpha) - m\,a_{\text{\tiny CM}_x} $$
\end{frame}


\begin{frame}
	\frametitle{Cálculo de aceleración angular y del CM (Tramo A)}
	Datos:
	$$ \begin{array}{c c c c c}
			F = \qty{20}{\newton} & \alpha = \qty{30}{\degree} & m = \qty{2}{\kilo\gram} & R= \qty{0.10}{\meter} & g =\qty[per-mode=fraction]{10}{\meter\per\second\squared}
		\end{array} $$
	Reemplazando en la ecuación
	$$ \frac{1}{2}\,m\,R\,\gamma = F\,\cos(\alpha) - m\,g\,\sin(\alpha) - m\,\gamma\,R  $$
	% $$ \frac{1}{2}\,m\,R\,\gamma = F\,\cos(\alpha) - m\,g\,\sin(\alpha) - m\,\gamma\,R  $$
	$$ \gamma = 2\,\frac{F\,\cos(\alpha) - m\,g\,\sin(\alpha)}{3\,m\,R} $$
	% $$ \gamma = 2\,\frac{F\,\cos(\alpha) - m\,g\,\sin(\alpha)}{3\,m\,R} \approx \qty[per-mode = fraction]{24.4}{\per\second\squared}
	% 	\implies a_{\text{\tiny CM}_x} \approx \qty[per-mode = fraction]{2.40}{\meter\per\second\squared} $$
	$$ \gamma = 2\cdot\frac{\qty{20}{\newton}\cdot\cos(\qty{30}{\degree}) - \qty{2}{\kilo\gram}\cdot\qty{10}{\meter\per\second\squared}\cdot\sin(\qty{30}{\degree})}{3\cdot\qty{2}{\kilo\gram}\cdot\qty{0.10}{\meter}} \approx \qty{24.40}{\per\second\squared} $$
	$$ \implies a_{\tcm} = \gamma\,R \approx \qty[per-mode = fraction]{2.44}{\meter\per\second\squared} $$

\end{frame}





\begin{frame}
	% \frametitle{Consideraciones Energéticas}
	% Como la fuerza de rozamiento es estática ya que solo actúa instantáneamente en el punto de contacto, no realiza trabajo.
	\frametitle{Trabajo - Contribuciones de las fuerzas}
	Analicemos qué fuerzas realizan trabajo
	\begin{columns}
		\column{0.5\textwidth}
		\begin{figure}
			\begin{tikzpicture}
				\planoInclinado[order=topdown, show floor]{0,2.5}{4,0}
				\def\R{1}

				\centroCirculoRampa{0.3}{\R}{C1}
				\dibujarCilindro[origen=C1]{1}
				\draw [dashed] (C1) -- (C1 |- 0,0) node[below]{$h_{\text{\tiny CM}_A}$};

				\centroCirculoRampa{0.85}{\R}{C2}
				\dibujarCilindro[origen=C2,rozamiento=1,ejes=true]{1}
				\draw [dashed] (C2) -- (C2 |- 0,0) node[below]{$h_{\text{\tiny CM}_C}$};
			\end{tikzpicture}
		\end{figure}
		\column{0.5\textwidth}
		% $$ \vec{F} = F\,(\cos(\alpha)\,\uveci + \sin(\alpha)\,\uvecj )\ \ \text{ y } \ \ m\,\vec{g} =  $$
		$$ m\,\vec{g} = m\,g\,(\sin(\alpha) \uveci + \cos(\alpha) \uvecj) $$
		$$ \vec{F} = F\,(\cos(\alpha)\,\uveci + \sin(\alpha)\,\uvecj )$$
		$$ \vec{N} = N\, \uvecj \ \ \ \text{ y } \ \ \ \vec{\fr}_e = -\fre\, \uveci$$
		\centering

		Y además
		$$ \mathrm{d}\vec{x} = \mathrm{d}x\,\uveci \implies \uveci \int_{x_0}^{x_f} \mathrm{d}x= \Delta x \, \uveci $$
	\end{columns}

	$$ W_{\vec{N}} = N\,\uvecj \cdot \Delta x\, \uveci  = 0$$
	$$ W_{\vec{\fre}_e} = -\fre\,\uveci \cdot \Delta x\, \uveci  = 0$$
	$$ W_{\vec{F}} =F\,(\cos(\alpha)\,\uveci + \sin(\alpha)\,\uvecj ) \cdot \Delta x \, \uveci $$
	$$ W_{m\vec{g}} =m\,g\,(\sin(\alpha)\,\uveci + \cos(\alpha)\,\uvecj ) \cdot \Delta x \, \uveci $$
\end{frame}

\begin{frame}
	\frametitle{Conservación de la Energía}
	Planteamos los teoremas de conservación de la energía
	$$ \Delta EC= EC_f - EC_o = WF_{NC} + WF_C $$
	$$ \Delta EP = - WF_C $$
	$$ \Delta EM = \Delta EC + \Delta EP = WF_{NC} $$

	Como el cuerpo parte del reposo
	$$ \Delta EC = EC_f $$

	Y la única fuerza no conservativa que realiza trabajo es $\vec{F}$
	$$ \Delta EM = EC_f + \Delta EP  = W_{\vec{F}}$$
\end{frame}

\begin{frame}
	\frametitle{Conservación de la Energía}
	\framesubtitle{Velocidad angular del CR}

	Determinamos la velocidad angular $\omega$ en el instante en el que deja de haber rozamiento.
	$$  \vec{v}_{\tcm} = \vec{v}_{\text{\tiny CIR}} + \vec{v}_{\tcm / \text{\tiny CIR}} = \vec{\omega} \times \vec{r}_{\tcm / \text{\tiny CIR}}$$
	$$\left| \begin{matrix}
			\uveci & \uvecj & \hat{k} \\
			0      & 0      & \omega  \\
			0      & -R     & 0
		\end{matrix}\right| $$
	$$ \vec{v}_\tcm = \omega\,R\, \uveci \implies \omega = \frac{v_\tcm}{R} $$

	Y como se demostró en~\eqref{eq:torque_conservation}, $\gamma = 0 \implies \omega = \text{cte.}$ durante todo el tramo $\overline{BC}$.


\end{frame}

\begin{frame}
	\frametitle{Conservación de la Energía}
	Tomando las definiciones de $EC$ y $EP$
	% $$ EC_f = \frac{1}{2}\,m\,v_{\tcm_f}^2 + \frac{1}{2}\, I_\tcm \,\omega^2 \ \ \text{y} \ \ \Delta EP = m\,g\,h_{\tcm_C} - m\,g\,h_{\tcm_A}  $$
	$$ EC_f = \frac{1}{2}\,m\,v_{\tcm_f}^2 + \frac{1}{2}\, I_\tcm \,\omega^2 $$
	$$\Delta EP = m\,g\,h_{\tcm_C} - m\,g\,h_{\tcm_A}  $$


	Y sabiendo que

	$$ \sin(\alpha) = \frac{h_{\tcm_C} - h_{\tcm_A}}{\Delta x} \implies \Delta x = \frac{h_{\tcm_C} - h_{\tcm_A}}{\sin(\alpha)}$$

	\vspace{1em}
	Desarrollamos la ecuación
	% $$ \frac{1}{2}\,m\,v_{\tcm_f}^2 + \frac{1}{4}\,m\,R^2\,\left( \frac{v_{\tcm_f}}{R}\right)^2  + m\,g\,\left(h_{\tcm_C} - h_{\tcm_A} \right) = \Delta x \, F \cos(\alpha)$$
	$$ \frac{1}{2}\,m\,v_{\tcm_f}^2 + \frac{1}{4}\,m\,R^2\,\left( \frac{v_{\tcm_f}}{R}\right)^2  + m\,g\,\left(h_{\tcm_C} - h_{\tcm_A} \right) = \frac{h_{\tcm_C} - h_{\tcm_A}}{\tan(\alpha)}F $$

	$$ \frac{3}{4}\,m\,v_{\tcm_f}^2 = (F\, \cot(\alpha) - m\,g)\,\left(h_{\tcm_C} - h_{\tcm_A} \right) $$

\end{frame}

\begin{frame}
	\frametitle{Conservación de la Energía}
	Teniendo en cuenta los datos del ejercicio
	$$ \begin{array}{c c c}
			F = \qty{20}{\newton} & m = \qty{2}{\kilo\gram} & g = \qty[per-mode=fraction]{10}{\meter\per\second\squared} \\ \alpha = \qty{30}{\degree} & h_{\tcm_A} = \qty{1}{\meter} & h_{\tcm_C} = \qty{4}{\meter}
		\end{array} $$

	$$ v_{\tcm} = \sqrt{\frac{F\, \cot(\alpha) - m\,g}{3\,m}\,4\,\left(h_{\tcm_C} - h_{\tcm_A} \right)} $$
	$$ v_{\tcm} = \sqrt{\frac{\qty{20}{\newton}\cdot \cot(\qty{30}{\degree}) - \qty{2}{\kilo\gram}\cdot\qty[per-mode=fraction]{10}{\meter\per\second\squared}}{3\cdot\qty{2}{\kilo\gram}}\cdot4\cdot\left(\qty{4}{\meter} - \qty{1}{\meter}\right)}$$

	$$ v_{\tcm_C} \approx \qty[per-mode = fraction]{5.41}{\meter\per\second} $$
	$$ \vec{v}_{\tcm_C} \approx \qty[per-mode = fraction]{5.41}{\meter\per\second} \, \uveci $$

\end{frame}

\end{document}
